\chapter{Conclusions and Future Work}
\label{ch:conclusions}

\section{Summary of findings}
\label{sec:concl-summary}

This thesis set out to break the circularity that undermines data-driven transit
planning and to build, in its place, a demand-driven pipeline that diagnoses unmet need,
generates candidate corridors, and transfers to other cities. The four research
questions can now be answered directly.

\paragraph{RQ1 -- Demand.} \emph{Yes.} Transit demand can be estimated at \ageb
resolution from Tier-1 data alone. The trip-generation and doubly-constrained gravity
model of \cref{sec:w1} produce a metropolitan demand surface using only census, economic,
and network variables, with no transit input, and calibration against the EOD~2022 survey
(\cref{sec:w2}) fixes the distance-decay parameter at $\betadecay=1.2005$---a value that
fits the observed flows better than the conventional prior. Demand is therefore never a
function of current supply, which is the precondition the rest of the framework depends on.

\paragraph{RQ2 -- Diagnosis.} \emph{Yes, and it validates out of sample.} The
coverage-gap index (\cref{sec:w3}) sets modeled demand against \textsc{gtfs}-measured
accessibility, yielding a dependent variable constructed independently of supply; the
supervised retrain learns it from place characteristics with no leakage
(\cref{tab:retrain}). Its strongest support is the Line~4 natural experiment: a corridor
the model never saw, later chosen for construction, shows \num{1.6} times the
metropolitan High-gap rate. This diagnostic is the thesis's firmest and most citable
contribution.

\paragraph{RQ3 -- Generation.} \emph{Partially, and now characterized.} Once corridors
are shaped as real paths (a minimum-spanning-tree diameter trunk on the road graph) and
judged by anchor-directness rather than endpoint detour (\cref{sec:w6}), the generator
produces at least one substantive, feasible, merit-passing corridor---\code{W6\_G02},
\SI{12.1}{\kilo\metre}, \SI{56}{\percent} High-gap, non-redundant, \num{73}rd-percentile
demand efficiency. It also surfaces lower-efficiency connectors alongside it, and it does
not reconstruct built rail lines---an expected outcome, since reconstruction is a weak
proxy for merit (\cref{sec:disc-reconstruction}). The generative layer is thus a genuine
but bounded contribution.

\paragraph{RQ4 -- Transferability.} \emph{Yes, end-to-end.} The full W1--W8 pipeline runs
without bespoke re-engineering on two contrasting metropolitan areas---Toluca (large,
sixteen municipalities) and Aguascalientes (compact, three)---re-deriving demand,
diagnosis, weights, corridors, and audit from each city's own Tier-1 and \textsc{gtfs}
data (\cref{ch:transferability}). The diagnostic differentiates the three metros
sensibly by unmet-demand share, and each city's weighting re-derives from its own data
rather than inheriting \zmg's.

\section{Practical recommendations}
\label{sec:concl-recommendations}

For a planning agency, the framework is most useful as a \emph{diagnostic screen} rather
than an autopilot. Three recommendations follow. First, treat the coverage-gap surface
as the primary product: it ranks where unmet demand is concentrated at a resolution finer
than any single line, and it can be recomputed cheaply as census and \textsc{gtfs} data
are updated. Second, use the route audit (\cref{sec:w7}) to flag low-demand, indirect,
and redundant existing routes for review---in the transfer cities it surfaced substantial
parallel-service redundancy that is a candidate for rationalization. Third, read the
generated corridors as \emph{hypotheses to be studied}, not as final alignments: the
merit-passing corridor is a defensible starting point for a feasibility study, but the
generator's characterized limitations mean its output should inform, not replace,
engineering and community process.

\section{Future work}
\label{sec:concl-future}

Four directions would most strengthen the framework. First, \emph{local calibration for
the transfer cities}: Toluca and Aguascalientes currently inherit \zmg's distance-decay
prior, and a survey-calibrated $\betadecay$ per city would tighten their demand surfaces.
Second, a \emph{network-distance gravity model}: replacing Euclidean impedance with
travel-time impedance on the road graph would improve absolute demand magnitudes. Third,
\emph{transferring the supervised layer}: the coverage-gap retrain was run only for \zmg,
and applying it to the transfer cities would test whether the learned relationship between
place characteristics and unmet demand is itself portable. Fourth, a \emph{hybrid corridor
shaper} that switches between diameter-trunk and travelling-salesman routing where the
latter remains feasible, which the exploratory analysis suggests would marginally improve
coverage without sacrificing feasibility. Beyond these, the clearest external validation
would come from observed ridership on any corridor that is eventually built along a
generated alignment---the counterfactual that, by construction, no purely ex-ante study
can supply.
